% !TITLE 魏晋
% !SUMMARY 文学的自觉时代
% !ORDER 0

\begin{intro}
魏晋是一个乱世，却也是中国文学最富创造力的时代之一。

东汉末年天下大乱，董卓进京、黄巾起义、群雄割据，数百年统一帝国分崩离析。但正是在这样的动荡中，文学第一次获得了独立的地位——不再是经学的附庸，不再是政教的工具，而是个人情感与生命意识的直接表达。鲁迅说魏晋是“文学的自觉时代”，这个判断至今仍然成立。

建安文学是魏晋的开端。曹操、曹丕、曹植父子三人，既是政治军事领袖，又是文坛巨擘。他们笔下的诗，没有了汉儒的迂腐说教，有的是直面乱世的慷慨悲凉。“白骨露于野，千里无鸡鸣”是曹操眼中的现实；“秋风萧瑟天气凉，草木摇落露为霜”是曹丕心中的感伤；“捐躯赴国难，视死忽如归”是曹植胸中的热血。这就是“建安风骨”——直面现实的勇气与深沉的生命意识。

到了正始年间，政治环境日趋险恶，司马氏专权，名士动辄得咎。嵇康、阮籍等人选择了另一种姿态：纵酒放达、谈玄论道，把对现实的失望转化为对个体精神世界的探索。阮籍的《咏怀诗》八十二首，晦涩曲折，却道出了一个知识分子在黑暗时代的全部焦虑。

陶渊明是魏晋文学的终点，也是新的起点。他不为五斗米折腰，归隐田园，在日常生活里发现了诗意。“采菊东篱下，悠然见南山”——这不是逃避，而是一种主动选择的生活方式。从建安的风骨到正始的玄思，再到陶渊明的自然，魏晋文学完成了从外在世界到内在精神的深度转向。
\end{intro}
